\section{文献注记：谁补了哪一块}\label{app:notes}

本讲义的写法是"原文为骨、教材为肉、我们补缝"。此附录说明每块肉来自哪本书、哪一节，
以及哪些地方我们查过原文页图确认过。完整的文献表（73 条）见工程目录
\texttt{materials/文献表.md}，本附录只列实际用到的。

\subsection{五本常用书的定位}
\begin{center}
\begin{tabular}{@{}p{2.6cm}p{2.0cm}p{7.6cm}@{}}
\toprule
书 & 定位 & 最适合补什么 \\
\midrule
河野俊丈《场论与拓扑学》（岩波 1998；中译 126 页） & 入口 & 仿射代数与 WZW（1.1--1.3）、共形块与融合规则（1.4）、KZ（1.5、2.1）、Jones 与三维不变量（2.2--2.3）、CS 与曲面联络空间（2.5）、Verlinde 等式（命题 1.22, p.37）、Dehn 手术（定理 2.16, p.60） \\
上野健爾《Conformal Field Theory with Gauge Symmetry》 & 细节手册 & 顶点代数、KZ 方程、单值性、共形块的形式定义（\S3.4、\S6.3） \\
Moore《Chern--Simons 理论导论》(TASI 2019) & 有疑查 & 相空间量子化与 Hamilton 约化（\S2.2.6）、曲面平坦联络（\S2.2.7）、环面量子化（\S2.2.8--12）、标架反常（\S2.2.14--17）、模空间量子化与代数几何（\S3.4） \\
Schottenloher《共形场论的数学导引》 & 数学侧 & Verlinde 公式的数学地位（第 11 章，书 p.172--189）：11.2 公式与证明史，11.3 融合规则，11.4 Verlinde 代数 \\
李思《量子力学与几何》 & 几何量子化 & 预量子化（5.4.3）、极化与量子态（5.4.4）、BKS、半形式（注：李思不讲半形式修正） \\
\bottomrule
\end{tabular}
\end{center}

\subsection{逐章对照}
\begin{center}
\begin{tabular}{@{}p{1.4cm}p{5.0cm}p{5.6cm}@{}}
\toprule
本讲义 & 补的主要来源 & 备注 \\
\midrule
第 2 讲 & 原文 \S1；河野 2.5 & 归一化换算见附录 A \\
第 3 讲 & 原文 \S3；Moore \S2.2.6--7；Atiyah--Bott & 辛商的三步是我们按定义重写的 \\
第 4 讲 & 李思 5.4.3--5.4.4；Moore \S3.4 & 半形式修正本文未做，$k+h^\vee$ 取自 (2.16)(2.17) 与后世解释 \\
第 5 讲 & 原文 \S3.2；河野 2.5 & "射影平坦"的原句 p.370 \\
第 6 讲 & 原文 \S3.3--3.4；Moore 3.3.2；BWB & BWB 的证明只给思路（Weyl 特征标/指标定理） \\
第 7 讲 & 河野 1.1--1.6；上野 \S3.4/\S6.3；KZ 原文 & 原文只引用，不展开 \\
第 8 讲 & 原文 \S4.1--4.5；河野 2.3--2.4；Moore \S4 & 手术公式的推导我们按"两种算法对比"重写 \\
第 9 讲 & 原文 \S4.5；河野 1.22/2.16；Schottenloher 第 11 章 & 十步链的编号是我们加的 \\
第 10 讲 & 原文 \S5；Brown--Henneaux；Maldacena & 1988 年无全息，须标年份 \\
第 11 讲 & Khovanov 2000；Witten 2011；Kapustin--Witten 2007 & 全部标"后世" \\
第 12 讲 & 原文 \S5 末 & 边界清单 \\
\bottomrule
\end{tabular}
\end{center}

\subsection{哪些我们核过原文页图}
\begin{itemize}
  \item (1.1)--(1.6)、(3.1)--(3.6)、(4.28)、(4.36)、(4.41)、(4.51)：逐符号对过 \texttt{pages/page-*.png}。
  \item "引言没有'本文解决第二问'这句话"：核过原刊引言与中译本导语框（导语框是译者所加）。
  \item 中译本 40 页、原文 49 页（351--399）：页边注「原文 p.N」是唯一可靠定位。
  \item 三处已知译本笔误：§3.2 与 p.368 把 (3.3) 误引成 (3.2)；(2.32) 相位因子与 p.365 的 $\delta$ 差 2；
        (1.2) 的 $1/4$ 已判否（原刊照印，无 $1/4$）。详见附录 D。
\end{itemize}

\subsection{没有用到的}
本讲义不依赖 1980 年代之前的原始文献（Atiyah 的会议报告、Ray--Singer、Quillen 等）：
这些只在需要精确出处时引用，不作为学习入口。理由是用户 2026-11-13 的指示：
"原始文献太老了，没有必要；今天的书写得更好"。例外是 Witten 本文本身——它是核心，必须逐节读。
